% paragraph 1: urbanization and increasing traffic congestion
% paragraph 2: Classical ramp metering methods
% paragraph 3: learning based / intelligent based traffic control
% paragraph 4: Bridge to Tech with RL and Applications
% paragraph 5: RL
%
%
%
%
%
%
% source un2018urbanizationnews: "68% of the world population is projected to live in urban areas by 2050" and "Urbanization and population growth could add another 2.5 billion people to urban areas by 2050"
% --> 1
%
% source chin2002temporary: "Traffic congestion and its impacts significantly affect the nation's economic performance and the public's quality of life"
% "in most urban areas, travel demand routinely exceeds highway capacity during peak periods"
% "Events such as crashes, vehicle breakdowns, work zones, adverse weather, and sub-optimal signal timing cause temporary capacity losses."
% --> 1
%
% source papageorgiou1997alinea:
% ALINEA is a local feedback ramp-metering strategy
% "The reported results indicate easy application, flexibility, and high efficiency of ALINEA."
% --> 2
%
% source yu2012genetic
% "Ramp metering is one of the most common solutions for improving motorway traffic flows."
% "Fuzzy logic based ramp metering algorithms are considered to be an effective control measure to handle a complex nonlinear traffic system."
% --> 2
%
%
% source ghods2009adaptive,
% "Ramp metering is the most widely used type of freeway traffic control."
% "Ramp metering, speed limits, and route recommendation are recognized as the most effective ways to relieve freeway traffic congestion.
% -->2
%
%
% source liang2009pid,
% "BP neural network is used to adjust the PID parameters dynamically."
% kann man verwenden: "Researchers also explored neural-network-based adaptive ramp metering approaches to dynamically regulate freeway traffic density [X]."
% --> 2 / 3
%
% source liang2010single,
% "A single neuron based traffic density controller is designed."
% "This method can effectively deal with this class of control problem."
% "It has good dynamic and steady-state performance, and can achieve an almost perfect tracking performance."
% kann man verwenden: "Neural-network-based approaches were also explored to dynamically regulate freeway traffic density and improve traffic stability [X]."
% --> 3
%
% source liang2009feedback,
% "This system can eliminate traffic congestion and maintain traffic flow stability."
% --> 3
%
%
% source qi2008feedback,
% Ramp metering is the most widely used control measure."
% "Neuron adaptive control algorithms are applied to tune the rate of metering on-line in real time."
% "Compared to traditional feedback ramp metering, these methods have stronger robustness and better control precision.
% --> 3
%
%
% source hegyi2005model,
% The use of dynamic speed limits significantly reduces congestion and results in a lower total time spent."
% --> 2
%
%
% source of feng2011fuzzy,
% Based on fuzzy metering and neural network theories…"
% "A five-layer Fuzzy Neural Network (FNN) model…"
% "Ramp metering rate as output variable…"
% --> 3
%
% source of sdot_connected_vehicle
% https://www.transportation.gov/briefing-room/us-dot-advances-deployment-connected-vehicle-technology-prevent-hundreds-thousands
% --> übergang zu RL, absatz 4
%
%
%
% Application in Games: mnih2013atari, silver2016go, silver2017go
% Applciation in Robotis: kober2013survey, polydoros2017survey
% Application in power systems: xu2019joint, zhang2020restoring
% Application in optimal building operation: zhang2021grid, zhang2020edge
% --> 4 --> Rl has achieved strong performance in many complex decision making domains
%
% Application in traffic signal control systems:
% sources eltantawy2010agent, jin2019multi, prashanth2011reinforcement, eltantawy2013multiagent, aslani2017adaptive,,
%
% Application in autonomous vehicles:
% sources: sallab2017deep, hoel2018automated, chae2017autonomous, ye2019automated, yu2020distributed, zhou2019efficient
%
%
% Applciation in speed limit control:
% sources wu2018differential, wu2019esctc
%
% Applciations in dynamic pricing of road usage:
% sources pandey2018multiagent, pandey2020pricing,
%
%
% source rezaee2014decentralized
% Dynamic traffic control measures provide a set of cost effective congestion mitigation solutions, among which ramp metering (RM) is the most effective approach. --> 1/2
% "The control system is based on reinforcement learning (RL)."
% "The proposed coordinated RM algorithm … significantly outperformed approaches based on the well-known ALINEA RM algorithm."
% "Resulted in 50% reduction in total travel time when applied to the Gardiner Model."
% --> 5
%
%
% source fares2014freeway,
% they use reinforcement learning control agents to freeway ramp metering
% --> 5
%
% source belletti2018expert,
% achieves state of the art performance or even outperforms it
%
% source watkins1992qlearning
% math behind Q learning
%
% source ummery1994online,
% math behind Q learning
%
%
% source konda2000actorcritic
% actor critic (PPO and A3C)
%
% source chulman2017ppo
% PPO
%
% source mnih2016a3c
% A3C
%
% source horgan2018distributed
% Ape-X DQN
%
% source wu2017flow
% Flow
%
% source krajzewicz2012recent
% SUMO
%
% source liang2018rllib
% RLlib
%
% source krajzewicz2014sumo
% SUMO emission model
%
%
% ---
%
%
% Text:


With the increase of urbanization and global population growth, cities are expected to house an additional 2.5 billion people by 2050 \cite{un2018urbanizationnews}.
This rising number of people in cities will inevitably result in higher traffic demand in metropolitan areas. Generally, travel demand surpasses
road capacity on a regular basis, especially with poorly coordinated traffic signal phases, and the resulting congestion negatively impacts not only a nation's economy, but also the individual's quality of life \cite{chin2002temporary}.

Freeway bottlenecks, particularly on-ramp merging areas, represent one of the major causes of congestion in urban transportation
networks \cite{chin2002temporary}.
To address this issue, ramp metering strategies have been widely adopted to regulate the flow of vehicles entering freeways and maintain stable
traffic conditions, where the red-to-green cycles are driven by a control algorithm \cite{yu2012genetic, ghods2009adaptive}.
Classical approaches that have been introduced in the literature include fixed-time controllers, feedback-based methods such as ALINEA \cite{papageorgiou1997alinea},
fuzzy-control \cite{yu2012genetic, ghods2009adaptive}, neural-network-based approaches \cite{liang2009pid, liang2010single, liang2009feedback, qi2008feedback},
model-based predictive controllers \cite{hegyi2005model} as well as techniques combining multiple control paradigms \cite{feng2011fuzzy}.

However, freeway traffic systems remain highly dynamic, stochastic and nonlinear, making it difficult for traditional control methods to generalize across varying traffic conditions.
In recent years, reinforcement learning (RL) has demonstrated impressive performance in solving complex decision-making problems through direct interaction with an environment.
These domains range from playing games \cite{mnih2013atari, silver2016go, silver2017go} to robotics \cite{kober2013survey, polydoros2017survey} as well as fields like energy systems \cite{xu2019joint, zhang2020restoring}
and intelligent infrastructure management \cite{zhang2021grid, zhang2020edge}.
Encouraged by these successes, RL has increasingly been adopted in mobility applications such as traffic signal systems \cite{eltantawy2010agent, jin2019multi, prashanth2011reinforcement, eltantawy2013multiagent, aslani2017adaptive},
decision making in autonomous vehicles \cite{sallab2017deep, hoel2018automated, chae2017autonomous, ye2019automated, yu2020distributed, zhou2019efficient},
dynamic speed regulation \cite{wu2018differential, wu2019esctc}
and adaptive road pricing strategies \cite{pandey2018multiagent, pandey2020pricing}.

Recently, several studies have also investigated the use of RL for freeway ramp metering control \cite{fares2014freeway, belletti2018expert}. For instance, Rezaee et al. \cite{rezaee2014decentralized} proposed an
RL-based approach that outperformed the widely used ALINEA controller.

Building upon these developments, this work is a reproduction study of the paper by Hou et al. \cite{hou2021rampmeter}, in which three different RL algorithms are applied
to the problem of freeway ramp metering control and compared to three baseline controllers.

